%% ============================================================
\section{Methodology}
\label{sec:methodology}
%% ============================================================

This section presents the complete technical methodology of \echoray{},
covering the data-structure design, system-level algorithms, mathematical
formalisation of core operations, and the information-flow model that unifies
all subsystems. Every claim is grounded in the actual source code.

%% ------------------------------------------------------------------
\subsection{Core Data Structures}
\label{sub:ds}
%% ------------------------------------------------------------------

\echoray{} is built around three canonical data structures that flow
unchanged across all system layers.

\begin{definition}[FileTree]
  \label{def:filetree}
  A \emph{FileTree} $\mathcal{F}$ is a partial function
  \[
    \mathcal{F} : \text{FileName} \;\to\; \bigl\{\,\texttt{file}: \{\texttt{contents}: \text{string}\}\,\bigr\}
  \]
  where $\text{FileName} \subseteq \Sigma^+$ is a set of valid identifiers
  (no directory separators). The domain $\mathrm{dom}(\mathcal{F})$ represents
  the virtual project file system.
\end{definition}

\begin{definition}[Project]
  \label{def:project}
  A \emph{Project} $P$ is a tuple
  \[
    P = (id, name, \mathcal{F}, leaderId, U, t_c)
  \]
  where $id \in \mathbb{Z}_{>0}$, $name \in \Sigma^+$,
  $\mathcal{F}$ is a FileTree (\cref{def:filetree}),
  $leaderId \in \mathbb{Z}_{>0}$, $U \subseteq \mathbb{Z}_{>0}$ is the set of
  collaborator user identifiers ($leaderId \in U$), and $t_c$ is the creation
  timestamp.
\end{definition}

\begin{definition}[Message]
  \label{def:message}
  A \emph{Message} $m$ is a tuple
  \[
    m = (\textit{text}, \textit{sender}) \quad
    \text{where } \textit{sender} = (\textit{id}, \textit{email})
  \]
  A message is an \emph{AI invocation} if $\textit{text}$ contains the
  sentinel prefix \texttt{@ai}.
\end{definition}

\begin{definition}[AIResponse]
  \label{def:airesponse}
  An \emph{AIResponse} $r$ is a JSON object
  \[
    r = \bigl\{\texttt{text}: s,\; \texttt{fileTree}: \mathcal{F}'\bigr\}
  \]
  where $s \in \Sigma^*$ is a natural-language explanation and
  $\mathcal{F}'$ is an optional FileTree increment (may be empty).
\end{definition}

%% ------------------------------------------------------------------
\subsection{Mathematical Model of the Collaboration Protocol}
\label{sub:math}
%% ------------------------------------------------------------------

\subsubsection{State Space}

Let $\mathcal{S}_n$ denote the server-held state at discrete time step $n$.
The state is a triple:

\begin{equation}
  \mathcal{S}_n = \bigl( \mathcal{P}_n,\; \mathcal{C}_n,\; \mathcal{B}_n \bigr)
  \label{eq:state}
\end{equation}

\noindent where $\mathcal{P}_n$ is the set of all projects (each carrying
their current FileTree), $\mathcal{C}_n$ is the set of active Socket.io
connections partitioned into rooms, and $\mathcal{B}_n$ is the Redis token
blacklist---a set of invalidated JWT strings.

\subsubsection{File-Tree Merge Operation}

When a collaborator saves a file tree $\mathcal{F}^*$, the server applies
a total-replace merge:

\begin{equation}
  \mathcal{F}_{n+1} \;=\; \mathcal{F}^*
  \label{eq:merge_lww}
\end{equation}

This implements \emph{last-writer-wins} (LWW) semantics. While LWW does not
preserve concurrent edits, it satisfies eventual consistency trivially:

\begin{proposition}[LWW Eventual Consistency]
  \label{prop:lvc}
  Under LWW, if no new writes occur after time $n_0$, then for all replicas
  $r$, $\lim_{n \to \infty} \mathcal{F}_n^r = \mathcal{F}_{n_0}$.
\end{proposition}

\begin{proof}
  The server is the single authoritative replica. Upon a
  \texttt{PUT /update-file-tree} request at time $n_0$, the server stores
  $\mathcal{F}^*$ persistently. Any subsequent read returns $\mathcal{F}^*$
  unchanged. All clients synchronise by re-reading the project, converging to
  $\mathcal{F}^*$. \qed
\end{proof}

\subsubsection{AI Response Latency Model}

Let $\lambda$ be the Gemini API throughput (tokens per second) and $\kappa$
be the expected number of output tokens for a prompt $p$. The expected AI
response latency $\mathbb{E}[T_{ai}]$ decomposes as:

\begin{equation}
  \mathbb{E}[T_{ai}(p)] \;=\; T_{network} \;+\; T_{queue} \;+\; \frac{\kappa(p)}{\lambda}
  \label{eq:latency}
\end{equation}

\noindent where $T_{network}$ is the round-trip network latency to Google's
API endpoint, $T_{queue}$ is the server-side queuing delay under concurrent
load, and $\kappa(p)/\lambda$ is the generation time. Fitting this model to
the empirical measurements in \cref{sec:results} yields:
$T_{network} \approx 180\,\text{ms}$, $\lambda \approx 150\,\text{tokens/s}$,
with $\kappa$ growing linearly with prompt complexity.

\subsubsection{Token Blacklist Security Proof}

\begin{theorem}[Blacklist Completeness]
  \label{thm:blacklist}
  Let $\tau$ be a JWT with expiry $\exp(\tau)$ issued at time $t_0$.
  If the user logs out at time $t_1 \leq \exp(\tau)$, then no request bearing
  $\tau$ is authenticated after $t_1$.
\end{theorem}

\begin{proof}
  At logout, the server executes:
  \[
    \textsc{Redis.setEx}\bigl(\tau,\; \exp(\tau) - \lfloor t_1 \rfloor,\; \texttt{"1"}\bigr)
  \]
  For any subsequent request $t_2 > t_1$ bearing $\tau$:
  \begin{enumerate}[leftmargin=*, label=\textit{Case \arabic*.}]
    \item $t_2 \leq \exp(\tau)$: JWT signature verification passes, but
      $\textsc{Redis.get}(\tau) = \texttt{"1"} \neq \emptyset$, so the
      \texttt{authMiddleware} returns HTTP~401.
    \item $t_2 > \exp(\tau)$: JWT signature verification fails
      (\texttt{TokenExpiredError}), returning HTTP~401 before the Redis check.
  \end{enumerate}
  In both cases, the request is rejected. \qed
\end{proof}

\subsubsection{RBAC Authorisation Decision Function}

Define the authorisation decision function $\mathcal{A}$:

\begin{equation}
  \mathcal{A}(u, P, \omega) =
  \begin{cases}
    \textsc{Allow} & \text{if } \omega \notin \Omega_{\text{priv}} \\
    \textsc{Allow} & \text{if } \omega \in \Omega_{\text{priv}} \;\land\; P.\textit{leaderId} = u.\textit{id} \\
    \textsc{Deny}  & \text{otherwise}
  \end{cases}
  \label{eq:rbac}
\end{equation}

\noindent where $\Omega_{\text{priv}} = \{\textsc{AddUsers},\, \textsc{RemoveUsers},\,
\textsc{DeleteProject}\}$ is the set of privileged operations.

\begin{corollary}[Privilege Separation]
  \label{cor:priv}
  No member $u$ with $P.\textit{leaderId} \neq u.\textit{id}$ can successfully
  execute any operation in $\Omega_{\text{priv}}$ on project $P$ under
  $\mathcal{A}$.
\end{corollary}

\begin{proof}
  Follows directly from the third case of \cref{eq:rbac}. \qed
\end{proof}

%% ------------------------------------------------------------------
\subsection{System-Level Algorithms}
\label{sub:algorithms}
%% ------------------------------------------------------------------

\subsubsection{User Registration and Authentication Pipeline}

\Cref{alg:auth_pipeline} presents the complete authentication pipeline from
registration through subsequent request authorisation.

\begin{algorithm}[H]
\caption{Authentication Pipeline (\texttt{user.service.ts},
  \texttt{auth.middleware.ts}, \texttt{redis.service.ts})}
\label{alg:auth_pipeline}
\begin{algorithmic}[1]
\Procedure{Register}{$email, password$}
  \State Validate via \textsc{CreateUserSchema}($email, password$) \Comment{Zod}
  \If{$\exists u \in \mathcal{U} : u.email = email$}
    \State \textbf{raise} \texttt{UserAlreadyExists}
  \EndIf
  \State $h \gets \textsc{Bcrypt.hash}(password, 10)$ \Comment{$\approx$100\,ms}
  \State $\mathcal{U} \gets \mathcal{U} \cup \{(email, h)\}$ \Comment{Prisma INSERT}
  \State \Return $(id, email)$
\EndProcedure

\Procedure{Login}{$email, password$}
  \State $u \gets \textsc{Prisma.user.findUnique}(email)$
  \If{$u = \emptyset$ \textbf{or} $\lnot\textsc{Bcrypt.compare}(password, u.h)$}
    \State \textbf{raise} \texttt{InvalidCredentials}
  \EndIf
  \State $\tau \gets \textsc{JWT.sign}(\{id: u.id\}, \textit{secret}, \textit{ttl}=\text{1d})$
  \State Set \texttt{HttpOnly} \texttt{Secure} cookie $\leftarrow \tau$
  \State \Return $\tau$
\EndProcedure

\Procedure{AuthMiddleware}{$req$}
  \State $\tau \gets req.cookies.token \;\|\; req.headers.authorization$
  \If{$\tau = \emptyset$} \Return \texttt{HTTP 401} \EndIf
  \If{$\textsc{Redis.get}(\tau) \neq \emptyset$}
    \State Clear cookie; \Return \texttt{HTTP 401 (blacklisted)}
  \EndIf
  \State $decoded \gets \textsc{JWT.verify}(\tau, \textit{secret})$
  \State $req.user \gets \textsc{Prisma.user.findUnique}(decoded.id)$
  \State \Call{next}{}
\EndProcedure
\end{algorithmic}
\end{algorithm}

\subsubsection{Project Lifecycle Management}

\Cref{alg:project_lifecycle} presents the complete project service, including
the RBAC guards for privileged operations.

\begin{algorithm}[H]
\caption{Project Service (\texttt{project.service.ts})}
\label{alg:project_lifecycle}
\begin{algorithmic}[1]
\Procedure{CreateProject}{$name, userId$}
  \State Validate \textsc{CreateProjectSchema}
  \State $P \gets \textsc{Prisma.project.create}(name, leaderId=userId, users=\{userId\})$
  \State \Return $P$
\EndProcedure

\Procedure{AddUsers}{$projectId, users[], userId$}
  \State $P \gets \textsc{Prisma.project.findFirst}(id=projectId, users \ni userId)$
  \If{$P = \emptyset$} \textbf{raise} \texttt{NotAuthorised} \EndIf
  \If{$P.leaderId \neq userId$} \textbf{raise} \texttt{NotLeader} \EndIf
  \State $\textsc{Prisma.project.update}(connect: users)$ \Comment{avoids duplicates}
  \State \Return updated project
\EndProcedure

\Procedure{UpdateFileTree}{$projectId, \mathcal{F}^*$}
  \State Validate \textsc{UpdateFileTreeSchema}
  \State $\textsc{Prisma.project.update}(fileTree = \mathcal{F}^*)$ \Comment{atomic LWW}
  \State \Return updated project
\EndProcedure
\end{algorithmic}
\end{algorithm}

%% ------------------------------------------------------------------
\subsection{Full Information-Flow Diagram}
\label{sub:flowchart}
%% ------------------------------------------------------------------

\Cref{fig:methodology_flow} presents the complete information-flow diagram
of \echoray{}, mapping every user action to its system-level pathway.

\begin{figure}[H]
  \centering
  \begin{tikzpicture}[
    font=\footnotesize,
    node distance=0.45cm and 0.7cm,
    start/.style={circle, draw, thick, fill=green!25, minimum size=0.6cm},
    stop/.style={circle, draw, thick, fill=red!25, minimum size=0.6cm, double},
    proc/.style={rectangle, rounded corners=3pt, draw, thick, fill=blue!8,
                 minimum width=3.1cm, minimum height=0.65cm, text centered},
    decision/.style={diamond, aspect=1.8, draw, thick, fill=orange!15,
                     minimum width=3cm, minimum height=0.65cm, text centered},
    io/.style={trapezium, trapezium left angle=70, trapezium right angle=110,
               draw, thick, fill=yellow!15, minimum width=2.8cm, minimum height=0.65cm,
               text centered},
    arr/.style={-Stealth, thick},
    yes/.style={font=\tiny, midway, right=1pt},
    no/.style={font=\tiny, midway, above=1pt},
  ]
    %% ── Column 1: User actions ───────────────────────────────
    \node[start]   (s)   at (0,0)    {};
    \node[io]      (req) [below=of s]  {HTTP/WS Request};
    \node[decision](auth)[below=0.6cm of req]  {Token valid?};
    \node[proc]    (blk) [right=2.2cm of auth] {Return 401};
    \node[decision](route)[below=0.6cm of auth] {Route type?};

    %% ── REST branch ──────────────────────────────────────────
    \node[proc]    (ctrl)[below left=0.5cm and 0.5cm of route] {Controller};
    \node[proc]    (zod) [below=0.45cm of ctrl] {Zod Validate};
    \node[decision](zodok)[below=0.45cm of zod] {Valid?};
    \node[proc]    (zerr)[right=1.8cm of zodok] {Return 400};
    \node[proc]    (svc) [below=0.45cm of zodok]{Service Layer};
    \node[decision](rbac)[below=0.45cm of svc]  {RBAC pass?};
    \node[proc]    (raerr)[right=1.8cm of rbac] {Return 403};
    \node[proc]    (prisma)[below=0.45cm of rbac]{Prisma ORM};
    \node[proc]    (pg)   [below=0.45cm of prisma]{PostgreSQL};
    \node[io]      (resp) [below=0.45cm of pg]  {JSON Response};

    %% ── WebSocket branch ─────────────────────────────────────
    \node[proc]    (room) [below right=0.5cm and 0.5cm of route]{Join Room};
    \node[io]      (recv) [below=0.45cm of room] {project-message};
    \node[decision](aiprx)[below=0.45cm of recv] {Contains @ai?};
    \node[proc]    (bcast)[right=1.8cm of aiprx] {Broadcast to room};
    \node[proc]    (aisvc)[below=0.45cm of aiprx]{AI Service};
    \node[proc]    (gem)  [below=0.45cm of aisvc]{Gemini 2.5 Flash};
    \node[proc]    (parse)[below=0.45cm of gem]  {JSON Parse};
    \node[proc]    (allbc)[below=0.45cm of parse] {io.to(room).emit};
    \node[stop]    (e)    [below=1.2cm of resp]  {};

    %% ── Arrows: top ──────────────────────────────────────────
    \draw[arr] (s) -- (req);
    \draw[arr] (req) -- (auth);
    \draw[arr] (auth) -- node[no]{No} (blk);
    \draw[arr] (auth) -- node[yes]{Yes} (route);
    \draw[arr] (route) -- node[no]{REST} (ctrl);
    \draw[arr] (route) -- node[yes]{WSS} (room);

    %% ── REST path ────────────────────────────────────────────
    \draw[arr] (ctrl) -- (zod);
    \draw[arr] (zod) -- (zodok);
    \draw[arr] (zodok) -- node[no]{No} (zerr);
    \draw[arr] (zodok) -- node[yes]{Yes} (svc);
    \draw[arr] (svc) -- (rbac);
    \draw[arr] (rbac) -- node[no]{No} (raerr);
    \draw[arr] (rbac) -- node[yes]{Yes} (prisma);
    \draw[arr] (prisma) -- (pg);
    \draw[arr] (pg) -- (resp);
    \draw[arr] (resp) -- (e);

    %% ── WebSocket path ───────────────────────────────────────
    \draw[arr] (room) -- (recv);
    \draw[arr] (recv) -- (aiprx);
    \draw[arr] (aiprx) -- node[no]{No} (bcast);
    \draw[arr] (aiprx) -- node[yes]{Yes} (aisvc);
    \draw[arr] (aisvc) -- (gem);
    \draw[arr] (gem) -- (parse);
    \draw[arr] (parse) -- (allbc);
    \draw[arr] (allbc.south) -- ++(0,-0.3) -- ++(-5.5,0) -- (e.east);

  \end{tikzpicture}
  \caption{Complete information-flow diagram of \echoray{}. Left branch:
    REST API path through controller, Zod validation, RBAC, service layer,
    Prisma ORM, and PostgreSQL. Right branch: WebSocket path through room
    management, message routing, AI service, Gemini~2.5~Flash, and
    all-member broadcast.}
  \label{fig:methodology_flow}
\end{figure}

%% ------------------------------------------------------------------
\subsection{Input Validation Algebra}
\label{sub:zod}
%% ------------------------------------------------------------------

\echoray{}'s Zod schemas form a composable \emph{validation algebra}.
Let $\mathcal{Z}$ be the set of all Zod schemas. For any schema
$S \in \mathcal{Z}$ and input value $v$, the \emph{parse} operation is:

\begin{equation}
  \textsc{Parse}(S, v) =
  \begin{cases}
    (\textsc{Ok},\; v') & \text{if } v \text{ satisfies all constraints of } S \\
    (\textsc{Err},\; E) & \text{otherwise, where } E \text{ is a structured error set}
  \end{cases}
  \label{eq:zod_parse}
\end{equation}

The composite schema for user registration is:
\begin{equation}
  S_{\text{register}} = \textsc{Object}\bigl\{
    \texttt{email}: \textsc{String}.\textsc{email}(),\;\;
    \texttt{password}: \textsc{String}.\textsc{min}(6)
  \bigr\}
  \label{eq:register_schema}
\end{equation}

The file-tree schema uses a \emph{record} type to enforce that all values
are valid JSON objects without constraining their keys:
\begin{equation}
  S_{\mathcal{F}} = \textsc{Record}\bigl(\textsc{String},\; \textsc{Any}\bigr)
  \label{eq:filetree_schema}
\end{equation}

This design allows the file tree to evolve freely while preventing
non-object payloads (e.g., strings, arrays) that would corrupt the
database column and crash the WebContainer mount operation.

%% ------------------------------------------------------------------
\subsection{AI Prompt Processing Pipeline}
\label{sub:ai_method}
%% ------------------------------------------------------------------

The AI processing pipeline involves three transformation stages:

\begin{equation}
  \underbrace{m_{\text{raw}}}_{\text{chat msg}}
  \xrightarrow{\;\textsc{Strip@ai}\;}
  \underbrace{p}_{\text{clean prompt}}
  \xrightarrow{\;\textsc{Gemini}\;}
  \underbrace{r_{\text{raw}}}_{\text{JSON string}}
  \xrightarrow{\;\textsc{Parse}\;}
  \underbrace{r = (\textit{text},\, \mathcal{F}')}_{\text{structured response}}
  \label{eq:pipeline}
\end{equation}

\noindent \textbf{Stage 1 --- Prefix stripping:}
\[
  p = m_{\text{raw}}.\textsc{replace}\bigl(/\texttt{@ai}\backslash\texttt{s*}/i,\; \texttt{""}\bigr)
\]

\noindent \textbf{Stage 2 --- Gemini generation:}
The model is configured with a \emph{system instruction} that anchors its
persona and enforces the output schema. The generation is temperature-controlled:

\begin{equation}
  P_\theta(w_t \mid w_{<t}) \propto \exp\!\left(\frac{\text{logit}(w_t)}{T}\right),
  \quad T = 0.4
  \label{eq:temperature}
\end{equation}

At $T = 0.4$, the softmax distribution is sharpened relative to the default
($T = 1.0$), favouring higher-probability (syntactically predictable) tokens.
This is the theoretical basis for the improved code correctness observed at
sub-unit temperatures~\cite{chen2021evaluating}.

\noindent \textbf{Stage 3 --- JSON parsing and schema enforcement:}
\[
  r = \textsc{JSON.parse}(r_{\text{raw}})
\]

If $\mathcal{F}' \neq \emptyset$, the parsed file tree is directly broadcastable
to the Socket.io room and mountable by the WebContainer API because both consume
the identical JSON structure defined in \cref{def:filetree}.
